% ============================================================
%  Chapter 6: R6 — Statistical Definition of Entropy
% ============================================================
\chapter{R6 — 熵的统计定义}
\label{ch:r6}

\section{总说：从微观到宏观的桥梁}

R6 是物理学中最深刻的桥梁之一。它连接了两个看似完全不同的世界：R4 的量子微观世界（单个粒子的态）和宏观热力学世界（温度、压力、热流）。

这条原理说：\textbf{熵——那个支配所有热力学过程的量——不是别的，就是一个系统在给定宏观约束下可以处于的微观状态数的对数}。

\begin{principlebox}[R6 — 熵的统计定义]
  \[
  S = k_B \ln W
  \]
  宏观系统的熵正比于系统在给定宏观约束（$E, V, N$ 等）下可及的微观状态总数的对数。

  \textbf{等概率先验假设}：孤立系统在平衡态时，所有可及的微观状态以相等概率出现（$P_i = 1/W$）。
  这是统计力学的唯一基本假设——所有其他结论（正则/巨正则分布、涨落-耗散定理等）都从此假设 + R4 + R6 推出。

  \textbf{热力学第二定律}：$\Delta S_{\text{total}} \ge 0$ 是 $S = k_B \ln W$ 和宇宙低熵初始条件的统计推论。
\end{principlebox}

% ── 变量定义表 ──
\section{变量定义}

\begin{table}[h]
\centering
\small
\begin{tabular}{p{2.2cm}p{1.6cm}p{2cm}p{2.2cm}p{3cm}p{3.5cm}}
\toprule
\textbf{变量} & \textbf{符号} & \textbf{类型} & \textbf{量纲} & \textbf{定义域} & \textbf{物理含义} \\
\midrule
熵 & $S$ & $\mathbb{R}_{\ge 0}$ & M$\cdot$L$^2\cdot$T$^{-2}\cdot\Theta^{-1}$ & $S \ge 0$ & 系统无序度的定量量度 \\
Boltzmann 常数 & $k_B$ & 常数 & M$\cdot$L$^2\cdot$T$^{-2}\cdot\Theta^{-1}$ & $1.381\times10^{-23}$ J/K & 能量-温度转换标度 \\
微观状态数 & $W$ & $\mathbb{N}^+$ & [1] & $W \ge 1$ & 可及量子微观状态总数 \\
温度 & $T$ & $\mathbb{R}_{\ge 0}$ & $\Theta$ & $T \ge 0$ & $1/T \equiv \partial S/\partial E$ \\
内能 & $E$ / $U$ & $\mathbb{R}$ & M$\cdot$L$^2\cdot$T$^{-2}$ & 由约束确定 & 系统总能量 \\
体积 & $V$ & $\mathbb{R}_{>0}$ & L$^3$ & $V > 0$ & 系统占据的空间 \\
粒子数 & $N$ & $\mathbb{N}$ & [1] & $N \ge 0$ & 微观粒子总数 \\
化学势 & $\mu$ & $\mathbb{R}$ & M$\cdot$L$^2\cdot$T$^{-2}$ & $\mu \equiv -T(\partial S/\partial N)_{E,V}$ & 增加一粒子所需能量 \\
微观态能量 & $E_i$ & $\mathbb{R}$ & M$\cdot$L$^2\cdot$T$^{-2}$ & $\hat{H}|i\rangle = E_i|i\rangle$ & 量子态能量本征值 \\
\bottomrule
\end{tabular}
\caption{R6 变量定义}
\end{table}

% ── 数学表达式 ──
\section{数学表达式}

\subsection{Boltzmann 熵公式}
\[
S(E, V, N) = k_B \ln W(E, V, N)
\]
$W(E,V,N)$ 为系统在总能量 $E$、体积 $V$、粒子数 $N$ 约束下的量子微观状态数。

\subsection{温度定义（由 $S = k_B \ln W$ 导出）}
\[
\frac{1}{T} \equiv \left(\frac{\partial S}{\partial E}\right)_{V, N}
\]

\subsection{热力学第二定律（统计形式）}
\[
\Delta S_{\text{total}} \ge 0 \quad \text{（孤立系统自发过程）}
\]
等号仅在可逆过程中成立。

\subsection{Boltzmann 分布（正则系综）}
\[
P_i = \frac{e^{-E_i/k_B T}}{Z}, \quad Z \equiv \sum_i e^{-E_i/k_B T}
\]

% ── 成立条件与失效边界 ──
\section{成立条件与失效边界}

\begin{limitbox}[R6 的成立条件与失效边界]
  \textbf{成立条件}：
  \begin{itemize}
    \item 宏观极限：$N \gg 1$（统计涨落 $\sim 1/\sqrt{N}$ 可忽略）
    \item 热平衡：$T$ 仅在平衡时有统一定义
    \item 遍历性：等概率假设仅当系统能遍历所有可及微观态时严格成立
  \end{itemize}

  \textbf{失效边界}：
  \begin{itemize}
    \item 小系统（$N \sim 1\text{--}100$）：统计涨落显著
    \item 非平衡过程：$T$ 无统一定义
    \item 自引力系统：热容为负；系综不等价
    \item 黑洞：Bekenstein-Hawking 熵 $S_{\text{BH}} = k_B c^3 A/(4G\hbar)$——正比于面积非体积
  \end{itemize}
\end{limitbox}

% ── 必要性 ──
\section{必要性 (Necessity)}

\begin{table}[h]
\centering
\small
\begin{tabular}{p{4cm}p{7cm}}
\toprule
\textbf{移除/弱化} & \textbf{后果} \\
\midrule
无 $S = k_B \ln W$ & 无微观-宏观桥梁；热力学仅为经验规律 \\
无熵增 & 无时间箭头；无热机效率限制 \\
无温度定义 $1/T \equiv \partial S/\partial E$ & 温度退化为经验标度 \\
无 $k_B$ & 熵-温度量纲独立；量子统计失去标度 \\
\bottomrule
\end{tabular}
\caption{R6 必要性}
\end{table}

% ── 充分性 ──
\section{充分性 (Sufficiency)}

\textbf{R6 + R4}（量子微观态）提供：
\begin{itemize}
  \item 热力学全部定律（第零/一/二/三）
  \item 理想气体状态方程 $PV = Nk_B T$
  \item 量子统计（Fermi-Dirac + Bose-Einstein）
  \item 化学平衡条件 $\mu_A = \mu_B$
  \item 相变理论
\end{itemize}
